\section{Characters and selection of a vacuum}
\label{ex-phase-selection}

The preceding section constructs a commutative algebra on the whole ground
space. Its characters provide a useful way of separating the vacuum state
into components. The argument below also explains why a one-dimensional
ground space is an almost-everywhere conclusion, rather than an assumption
about the unselected table.

\begin{proposition}[The representing probability]\label{ex-character-measure}
Suppose the algebra \(\mathcal A\) of
Proposition~\ref{ex-ground-commutative} is separable in operator norm.
Its space \(\Xi\) of unital complex star characters, with the topology
of pointwise convergence on \(\mathcal A\), is compact and metrizable.
There is a probability measure \(\mu\) on \(\Xi\), positive on every
nonempty open set, such that
\begin{equation}\label{ex-spectral-state}
 \int_\Xi\chi(A)\,d\mu(\chi)=\langle\Omega,A\Omega\rangle
 \quad(A\in\mathcal A).
\end{equation}
\end{proposition}
\begin{proof}
The algebra is a nonzero unital commutative \(C^*\)-algebra: it is a
norm-closed star subalgebra of \(B(\mathcal H_0)\), and
\(\|\Omega\|=1\). The commutative Gelfand theorem identifies it
isometrically, as a star algebra, with \(C(\Xi)\), by
\(A\mapsto\widehat A\), where \(\widehat A(\chi)=\chi(A)\)
\cite[Theorem 2.13]{Williams}. Its character space is compact since the
algebra is unital.

Choose a countable norm-dense sequence \((A_j)\) in \(\mathcal A\).
The continuous map \(\chi\mapsto(\chi(A_j))_j\) into
\(\mathbb C^{\mathbb N}\) is injective: equality on the dense sequence
implies equality everywhere by continuity of the characters. A continuous
injection of a compact space into a Hausdorff space is a homeomorphism onto
its image (images of closed sets are compact and hence closed).
The countable product is metrizable, for example by
\(\sum_j2^{-j-1}\min(1,|z_j-z'_j|)\), so \(\Xi\) is metrizable.

Transfer the functional \(A\mapsto\langle\Omega,A\Omega\rangle\)
to \(C(\Xi)\). It is positive: a nonnegative continuous function is
\(\overline g g\) for its continuous nonnegative square root \(g\),
and its inverse Gelfand transform has the form \(B^*B\). The functional
takes the value \(\|B\Omega\|^2\ge0\) there. The Riesz representation
theorem supplies a positive Borel measure representing it
\cite[Chapter 6, Theorem 32]{Sawyer}. Compactness gives
\(C_c(\Xi)=C(\Xi)\), so that theorem applies directly. At the constant
function one it gives \(\mu(\Xi)=\|\Omega\|^2=1\).

To prove full support, suppose a nonempty open set \(O\) has zero measure.
Choose \(x\in O\) and \(r>0\) with \(B(x,r)\subset O\). The function
\(f(y)=\max(0,r-\operatorname{dist}(x,y))\) is continuous, nonzero at
\(x\), and zero outside \(O\). Let \(A\) be its inverse Gelfand
transform. Then \eqref{ex-spectral-state} applied to \(A^*A\) gives
\(\|A\Omega\|^2=\int|f|^2\,d\mu=0\). Separation of \(\Omega\)
forces \(A=0\), contradicting \(f(x)=r\). Thus every such open set has
positive measure.
\end{proof}

For \(\chi\in\Xi\), define the whole-word values
\begin{equation}\label{ex-phase-word}
 S_\chi(w)=\chi(M(w))
 \quad\text{for distinct-point words }w.
\end{equation}
It is essential to apply the character to the operator for the entire
word. The compressed operator \(M(w_1w_2)\) need not equal
\(M(w_1)M(w_2)\) at finite separation.
Equations~\eqref{ex-spectral-state} and \eqref{ex-compression-values} give
the barycenter identity
\begin{equation}\label{ex-phase-barycenter}
 S(w)=\int_\Xi S_\chi(w)\,d\mu(\chi).
\end{equation}
For a fixed word the integrand is continuous, and hence integrable.

Some properties of \(S_\chi\) already follow at this level. Normalization
and translation invariance follow from \(M(\varnothing)=I\) and
\eqref{ex-compression-values}. It is also reflection positive. For
positive words \(v_i\) and coefficients \(c_i\), the Gram identity in
the preceding section gives
\[
 \sum_{i,j}\overline{c_i}c_jM((\theta v_i)v_j)
 =\left(\sum_i c_iX(v_i)\right)^*
   \left(\sum_j c_jX(v_j)\right)\ge0
 \quad\text{in }B(\mathcal H_0).
\]
The operator on the left lies in \(\mathcal A\). Positivity is the
same in a unital \(C^*\)-subalgebra as in the ambient operator algebra
(spectral permanence and continuous functional calculus give its nonnegative square root in
the subalgebra; see \cite[Theorems 2.17 and 2.19]{Williams}). Characters take positive elements to nonnegative
numbers. Applying \(\chi\) proves the desired Gram positivity.
Regularity and the remaining covariance and field identities still need
their own operator estimates; the barycenter formula alone does not
establish them.

\begin{lemma}[A common countable collection of words]\label{ex-phase-horizon}
Suppose the raw label spaces are finite-dimensional complex vector spaces
indexed by the nonnegative integers. Suppose that each \(S_\chi\) is
multilinear in its labels and continuous on the distinct-point
configuration space, and has the Hilbert realization of
Proposition~\ref{ex-gns}. There is one countable collection of positive
words whose vectors span a dense subspace in every \(\mathcal H_\chi\).
\end{lemma}
\begin{proof}
Choose a finite basis in every raw label space. Take all positive words
with rational coordinates and with basis-vector labels, for every arity
and every finite list of label grades; include the empty word. There are
countably many choices. Every word is a finite linear combination of
words with basis labels, by expanding each label in its basis. This identity holds for the Hilbert-space vectors: pair the difference with any word vector and use multilinearity of the correlation in the positive-word labels. All such inner products vanish, so density forces the difference to be zero.
For a fixed basis-labelled positive word, approximate its positions by
rational positions, remaining within the open set of positive
distinct-point configurations. Its vectors converge in every
\(\mathcal H_\chi\): the squared norm of their difference is the sum
of two diagonal Gram values minus twice the real part of the cross Gram
value, and continuity makes this tend to zero. The reflected and positive
halves remain strictly separated throughout this limit. The full word
span is dense by construction; thus the selected countable span is dense
as well. The positions and bases were chosen without reference to
\(\chi\), so the collection is common to every phase.
\end{proof}

\begin{theorem}[A one-dimensional ground space in almost every phase]
\label{ex-pure-ground}
Assume every \(S_\chi\) has the Euclidean Hilbert-space dynamics of
Theorem~\ref{ex-euclidean-semigroup}, with unit vacuum \(\Omega_\chi\).
Assume there is a common countable collection of positive words whose
vectors span densely in every phase Hilbert space. Then, for
\(\mu\)-almost every \(\chi\),
\[
 \mathcal H_{\chi,0}=\mathbb C\Omega_\chi.
\]
In particular there exists a phase with this property.
\end{theorem}
\begin{proof}
Fix a positive word \(v\). For integers \(n\ge0\), set
\[
 f_n(\chi)=\|E_{\chi,n}[v]_\chi\|^2
 =S_\chi((\theta\tau_nv)(\tau_nv)).
\]
Each is a continuous function of \(\chi\), by
\eqref{ex-phase-word}. Contractivity gives \(0\le f_n\le f_0\);
the latter is integrable on the compact probability space. The
large-time projection theorem gives
\(f_n(\chi)\to\|P_{\chi,0}[v]_\chi\|^2\).
Dominated convergence and the barycenter identity imply
\begin{equation}\label{ex-phase-ground-average}
 \int_\Xi\|P_{\chi,0}[v]_\chi\|^2\,d\mu
 =\lim_{n\to\infty}\|E_n[v]\|^2=\|P_0[v]\|^2.
\end{equation}
All limit functions here are measurable as pointwise limits of the
continuous \(f_n\).

Apply \eqref{ex-spectral-state} to \(M(v)^*M(v)\). Since characters
preserve products and adjoints and \(M(v)\Omega=P_0[v]\),
\[
 \int_\Xi|S_\chi(v)|^2\,d\mu=\|P_0[v]\|^2.
\]
Moreover \(\langle\Omega_\chi,P_{\chi,0}[v]_\chi\rangle=S_\chi(v)\),
because the projection fixes the vacuum. Orthogonal decomposition along
the unit vector \(\Omega_\chi\) gives
\[
 \|P_{\chi,0}[v]_\chi-S_\chi(v)\Omega_\chi\|^2
 =\|P_{\chi,0}[v]_\chi\|^2-|S_\chi(v)|^2.
\]
The nonnegative right side is integrable and, by the two preceding
identities, has integral zero. It is therefore zero almost everywhere.

Apply this conclusion to the common countable collection and intersect
its full-measure sets. For every \(\chi\) in the intersection, the
projection of its dense word span lies in \(\mathbb C\Omega_\chi\).
This one-dimensional subspace is closed. Continuity of the projection
therefore implies \(\operatorname{ran}P_{\chi,0}\subset
\mathbb C\Omega_\chi\). The reverse inclusion holds because the vacuum
is fixed. Finally a set of measure one in a probability space is nonempty.
\end{proof}

This selection argument requires the phase Hilbert spaces and their
contraction semigroups before it can be applied. It does not construct
them from a probability decomposition alone. Retention of the raw-label
null spaces, field derivatives, stress tensor and relativistic
reconstruction remain additional obligations for the chosen phase.
